% !PRINTED XIII
\weberpage{XIII}
\inhaltpar{119}{Convergenz der Näherungsbrüche}{369}
\inhaltpar{120}{Aequivalente Zahlen}{371}
\inhaltpar{121}{Entwickelung äquivalenter Zahlen in Kettenbrüche}{374}
\inhaltpar{122}{Quadratische Irrationalzahlen}{377}
\inhaltpar{123}{Reducirte Zahlen mit negativer Discriminante}{379}
\inhaltpar{124}{Reducirte Zahlen mit positiver Discriminante}{383}
\inhaltpar{125}{Entwickelung reeller quadratischer Irrationalzahlen in Kettenbrüche}{386}
\inhaltpar{126}{Beispiele}{392}
\inhaltpar{127}{Die \gesperrt{Pell}'sche Gleichung}{395}
\inhaltpar{128}{Ableitung aller Lösungen der \gesperrt{Pell}'schen Gleichung aus der kleinsten positiven}{399}
\inhaltpar{129}{Genäherte Berechnung der reellen Wurzeln einer numerischen Gleichung durch Kettenbrüche}{401}
\inhaltpar{130}{Rationale Wurzeln ganzzahliger Gleichungen. Reducible Gleichungen}{403}
\inhaltabschnitt{Zwölfter Abschnitt.}{Theorie der Einheitswurzeln.}
\inhaltpar{131}{Die Einheitswurzeln}{408}
\inhaltpar{132}{Primitive Einheitswurzeln}{411}
\inhaltpar{133}{Gleichungen für die primitiven Einheitswurzeln $n^{\text{ten}}$ Grades}{414}
\inhaltpar{134}{Irreducibilität}{417}
\inhaltpar{135}{Die Discriminante der Kreistheilungsgleichung}{421}
\inhaltpar{136}{Primitive Congruenzwurzeln}{425}
\inhaltpar{137}{Multiplikation und Theilung der trigonometrischen Functionen}{432}
\inhaltpar{138}{Vorzeichenbestimmung. Quadratische Reste}{439}
\inhaltbuch{Drittes Buch.}{Algebraische Grössen.}
\inhaltabschnitt{Dreizehnter Abschnitt.}{Die Galois'sche Theorie.}
\inhaltpar{139}{Der Körperbegriff}{449}
\inhaltpar{140}{Adjunction}{451}
\inhaltpar{141}{Functionen in einem Körper}{452}
\inhaltpar{142}{Algebraische Körper}{455}
\inhaltpar{143}{Gleichzeitige Adjunction mehrerer algebraischer Grössen}{457}
\inhaltpar{144}{Primitive und imprimitive Körper}{460}
\inhaltpar{145}{Normalkörper. \gesperrt{Galois}'sche Resolvente}{464}
\inhaltpar{146}{Die Substitutionen eines Normalkörpers}{467}
\inhaltpar{147}{Zusammensetzung von Substitutionen}{470}
\inhaltpar{148}{Permutationsgruppen}{472}
\inhaltpar{149}{Galois'sche Gruppe}{476}
